%% GENERATED by tools/gen_latex.py -- edit content/ch14.py instead.
%% Build:  latexmk -pdf ch14.tex     (or: pdflatex ch14.tex, twice)
\documentclass[aspectratio=169,11pt,t]{beamer}
\usepackage{itbpro}

\coursetitle{Text Classification}
\coursesubtitle{Tokenization, and the one question every NLP architecture answers differently: is a sentence a set of words, or a sequence of them?}
\coursekicker{Chapter 14}
\coursesource{Chollet \& Watson, Deep Learning with Python 3e --- chapter 14}
\courseduration{3 hours (2 sessions)}
\coursepresenter{Rahman Indra Kesuma, S.Kom., M.Cs.}{Teaching Assistant}
\coursebrand{AI for Professional \textperiodcentered\ ITB \textperiodcentered\ Ch 14}

\title{Text Classification}
\author{Rahman Indra Kesuma, S.Kom., M.Cs.}
\date{}

\begin{document}
\covertitle

\framekicker{Learning outcomes}
\begin{frame}[shrink=25]{Session Objectives}
\leadpar{By the end of this session, participants can:}
\begin{isteps}
\item Explain why natural language is \textbf{messy in a way machine languages are not}, and why rule-based NLP failed.
\item Compare \textbf{character, word, and subword} tokenization, and say what byte-pair encoding optimises.
\item State the \textbf{set versus sequence} question and place bag-of-words, RNNs, and Transformers against it.
\item Build a \textbf{bag-of-words} classifier with \inlinecode{TextVectorization}, and extend it with \textbf{bigrams}.
\item Explain what a \textbf{word embedding} is, and why one-hot encoding makes a false assumption.
\item \textbf{Pretrain} an embedding with an unsupervised task (CBOW), and use it for classification.
\end{isteps}

\vskip 6pt
\vskip 4pt
\reslink{Book}{Chapter 14 --- full text}{https://deeplearningwithpython.io/chapters/chapter14\_text-classification}
\reslink{Notebook}{01\_tokenizers.ipynb}{../../notebooks/ch14/01\_tokenizers.ipynb}
\reslink{Notebook}{02\_bag\_of\_words\_and\_bigrams.ipynb}{../../notebooks/ch14/02\_bag\_of\_words\_and\_bigrams.ipynb}
\reslink{Notebook}{03\_sequence\_models.ipynb}{../../notebooks/ch14/03\_sequence\_models.ipynb}
\reslink{Notebook}{04\_pretraining\_an\_embedding.ipynb}{../../notebooks/ch14/04\_pretraining\_an\_embedding.ipynb}
\reslink{Repo}{Author's official notebook}{https://github.com/fchollet/deep-learning-with-python-notebooks/blob/master/chapter14\_text-classification.ipynb}
\vskip 2pt
\end{frame}

\framekicker{Section 14.1}
\begin{frame}[shrink=25]{Why "natural" language is a different kind of object}
\begin{minipage}[t]{0.485\textwidth}
{\fontsize{9.4}{12.6}\selectfont\color{fg2}\textbf{Machine languages were designed.} An engineer wrote down formal rules first; people used the language once the rule set was complete.\par}\vskip 4pt

\end{minipage}%
\hfill
\begin{minipage}[t]{0.485\textwidth}
{\fontsize{9.4}{12.6}\selectfont\color{fg2}\textbf{Human language is the reverse.} Usage comes first; rules arise later, are formalised after the fact, and are \hilite{often ignored or broken by its users}\par}\vskip 4pt

\end{minipage}%
\begin{iband}[signal]
So machine-readable language is highly structured and rigorous, while natural language is \textbf{messy --- ambiguous, chaotic, sprawling, and constantly in flux}.
\end{iband}
\end{frame}

\framekicker{Section 14.1}
\begin{frame}[shrink=25]{The 1954 demo, and the decade it bought}
{\fontsize{9.4}{12.6}\selectfont\color{fg2}Researchers at IBM and Georgetown showed a system translating Russian into English using \textbf{six hardcoded grammar rules} and a lookup table of a couple of hundred entries, on \textbf{60 handpicked sentences}.\par}\vskip 4pt
\begin{iband}[amber]
The goal was to attract excitement and funding, and in that sense \textbf{it was a huge success}. The authors claimed translation would be solved within five years. \hilite{Funding poured in for the better part of a decade.}
\end{iband}
\end{frame}
\note{Worth connecting back to chapter 1: the same pattern of a compelling demo, an over-confident timeline, and a funding collapse.}

\framekicker{Section 14.1}
\begin{frame}[shrink=25]{\ldots{}and why it did not generalise}
\begin{ibullets}
\item \textbf{Words change meaning dramatically depending on context.}
\item \textbf{Every grammar rule needed countless exceptions.}
\item Shining on a few handpicked examples was simple; building a robust system that could compete with human translators \hilite{was another matter}.
\end{ibullets}
{\fontsize{9.4}{12.6}\selectfont\color{fg2}An influential US report a decade later picked apart the lack of progress, and the funding dried up. Yet handcrafted rules \textbf{remained dominant well into the 1990s} --- because there was no viable alternative.\par}\vskip 4pt
\end{frame}

\framekicker{Section 14.1}
\begin{frame}[shrink=25]{The moment it became machine learning}
\begin{center}
\includegraphics[width=0.92\linewidth,height=0.60\textheight,keepaspectratio]{figs/ch14-history.tex.pdf}
\end{center}
\vskip 3pt{\color{ink3}\fontsize{7.4}{9.4}\selectfont Faster computers and more data in the late 1980s changed what was possible.\par}
\begin{iquote}
\quotetext{When you find yourself building systems that are big piles of ad hoc rules \ldots{} you're likely to start asking, "Could I use a corpus of data to automate the process of finding these rules?" And just like that, you've graduated to doing machine learning.}\quotecite{Chollet \& Watson, section 14.1}
\end{iquote}
\end{frame}

\coursesection{01}{Preparing text data}{Tokenization: from characters to byte-pair encoding.}

\framekicker{Section 14.2.1 $\cdot$ listing 14.1}
\begin{frame}[shrink=25]{The simplest tokenizer, written out}
{\fontsize{9.4}{12.6}\selectfont\color{fg2}Every tokenizer does the same three things: standardise, split, index. Here they are at character level, which makes the structure easy to see.\par}\vskip 4pt
\icode{python}{listing 14.1 --- a character-level tokenizer}{listings/ch14-01.py}
{\fontsize{9.4}{12.6}\selectfont\color{fg2}The \inlinecode{[UNK]} fallback is the important detail: \textbf{anything outside the vocabulary becomes one shared token}, and how often that happens is a property of the tokenizer's design.\par}\vskip 4pt
\end{frame}

\framekicker{Section 14.2.1}
\begin{frame}[shrink=25]{Building a vocabulary, and capping it}
{\fontsize{9.4}{12.6}\selectfont\color{fg2}A tokenizer needs a vocabulary. Rather than mapping every distinct token to an index, it is capped at the most common ones --- which is a modelling decision, not a technicality.\par}\vskip 4pt
\icode{python}{computing a capped vocabulary}{listings/ch14-02.py}
\begin{iband}[signal]
Capping the vocabulary is the same \textbf{feature selection} move as chapter 5's \inlinecode{num\_words=10000}: rare tokens carry little signal and \hilite{invite the spurious correlations chapter 5 warned about}
\end{iband}
\end{frame}

\framekicker{Section 14.2.1}
\begin{frame}[shrink=25]{Word level: shorter sequences, a harder vocabulary problem}
{\fontsize{9.4}{12.6}\selectfont\color{fg2}Splitting on whitespace instead of on characters gives far shorter sequences --- and immediately raises two questions the character tokenizer never had.\par}\vskip 4pt
\icode{python}{a word-level split}{listings/ch14-03.py}
\begin{ibullets}
\item \textbf{How large should the vocabulary be?} English has hundreds of thousands of word forms.
\item \textbf{What happens to a word you have never seen?} It becomes \inlinecode{[UNK]}, and \hilite{all its meaning is gone}.
\end{ibullets}
\end{frame}

\framekicker{Section 14.2.1 -- 14.2.2}
\begin{frame}[shrink=25]{Three levels, and the trade-off between them}
\begin{center}
\includegraphics[width=0.92\linewidth,height=0.60\textheight,keepaspectratio]{figs/ch14-tokenizers.tex.pdf}
\end{center}
\vskip 3pt{\color{ink3}\fontsize{7.4}{9.4}\selectfont Subword tokenization aims at the best of both.\par}
\end{frame}

\framekicker{Section 14.2.2}
\begin{frame}[shrink=25]{Tokenization as compression}
{\fontsize{9.4}{12.6}\selectfont\color{fg2}The book reframes the search for a good tokenizer as \textbf{the hunt for a good compression of the input}.\par}\vskip 4pt
\begin{ibullets}
\item \textbf{Shorter tokens} compress the overall length of each example.
\item \textbf{A smaller vocabulary} reduces the bytes needed to represent each token.
\item Achieve both and you feed the model \hilite{short, information-rich sequences}.
\end{ibullets}
\begin{iband}[signal]
That analogy turned out to be powerful: one of the most effective tricks of the last decade was \textbf{repurposing a 1990s lossless compression algorithm --- byte-pair encoding --- for tokenization}. It is used by \hilite{ChatGPT and many other models to this day}
\end{iband}
\end{frame}

\framekicker{Section 14.2.2}
\begin{frame}[shrink=25]{How byte-pair encoding builds its vocabulary}
{\fontsize{9.4}{12.6}\selectfont\color{fg2}Start from characters, then repeatedly merge the most frequent adjacent pair into a new token.\par}\vskip 4pt
\icode{python}{the BPE idea, in outline}{listings/ch14-04.py}
\begin{iband}[signal]
Common words end up as \textbf{single tokens}; rare words decompose into \textbf{familiar pieces}. Nothing is ever \inlinecode{[UNK]}, because \hilite{the character level is always available as a fallback}
\end{iband}
\end{frame}

\coursesection{02}{Sets versus sequences}{The pivotal question from which every NLP architecture springs.}

\framekicker{Section 14.3}
\begin{frame}[shrink=25]{Representing a token is easy. Representing order is not.}
{\fontsize{9.4}{12.6}\selectfont\color{fg2}Individual tokens are \textbf{categorical features}, and we know how to handle those. The problematic question is how to encode \textbf{the ordering} of tokens.\par}\vskip 4pt
\begin{iband}[signal]
Unlike the steps of a timeseries, \textbf{words in a sentence have no natural, canonical order}. Different languages order similar words very differently; within one language you can usually say the same thing by reshuffling. \hilite{Order clearly matters, but its relationship to meaning is not straightforward.}
\end{iband}
\end{frame}

\framekicker{Section 14.3}
\begin{frame}[shrink=25]{Three answers, three families of architecture}
\begin{center}
\includegraphics[width=0.92\linewidth,height=0.60\textheight,keepaspectratio]{figs/ch14-setseq.tex.pdf}
\end{center}
\vskip 3pt{\color{ink3}\fontsize{7.4}{9.4}\selectfont RNNs and Transformers both take order into account, so both are called sequence models.\par}
{\fontsize{9.4}{12.6}\selectfont\color{fg2}Historically most early NLP was bag-of-words. Interest in sequence models only rose in \textbf{2015}, with the rebirth of RNNs. \textbf{Both approaches remain relevant today.}\par}\vskip 4pt
\end{frame}

\framekicker{Section 14.3.1}
\begin{frame}[shrink=25]{The benchmark: raw IMDB this time}
{\fontsize{9.4}{12.6}\selectfont\color{fg2}Chapters 4 and 5 used a \textbf{prevectorised} IMDB. Here the raw text is processed --- as you would with any new text problem.\par}\vskip 4pt
\icode{python}{listing 14.8 --- downloading raw IMDB}{listings/ch14-05.py}
\ioutput{listings/ch14-06.txt}
{\fontsize{9.4}{12.6}\selectfont\color{fg2}Note \inlinecode{train/unsup} --- 25,000 \textbf{unlabelled} reviews. They are ignored at first and become essential in section 14.5.4.\par}\vskip 4pt
\end{frame}

\coursesection{03}{Set models}{Throw order away, and see how far that gets you.}

\framekicker{Section 14.4}
\begin{frame}[shrink=25]{Bag of words, in one picture}
\begin{center}
\includegraphics[width=0.92\linewidth,height=0.60\textheight,keepaspectratio]{figs/ch14-bow.tex.pdf}
\end{center}
\vskip 3pt{\color{ink3}\fontsize{7.4}{9.4}\selectfont Tokenize, discard order, multi-hot encode.\par}
{\fontsize{9.4}{12.6}\selectfont\color{fg2}The idea is simply to \textbf{assign a weight to every word}. \textit{terrible} probably indicates a bad review; \textit{riveting} probably indicates a good one.\par}\vskip 4pt
\end{frame}

\framekicker{Section 14.4.1 $\cdot$ listing 14.12}
\begin{frame}[shrink=25]{TextVectorization does all of it}
{\fontsize{9.4}{12.6}\selectfont\color{fg2}Keras has a preprocessing layer for this. Like several such layers it has an \textbf{\inlinecode{adapt()}} method that learns state --- here, the vocabulary --- by iterating over a dataset.\par}\vskip 4pt
\icode{python}{listing 14.12 --- bag-of-words encoding}{listings/ch14-07.py}
\begin{iband}[signal]
\textbf{20,000 words is a good starting point} for text classification. Note \inlinecode{adapt()} is given the reviews \textbf{without labels} --- it is learning vocabulary, \hilite{not learning the task}
\end{iband}
\end{frame}

\framekicker{Section 14.4.1}
\begin{frame}[shrink=25]{The model on top is tiny}
{\fontsize{9.4}{12.6}\selectfont\color{fg2}Once the text is a 20,000-long multi-hot vector, the classifier is the same kind of small dense network as chapter 4.\par}\vskip 4pt
\icode{python}{the bag-of-words classifier}{listings/ch14-08.py}
{\fontsize{9.4}{12.6}\selectfont\color{fg2}It performs respectably --- which is the point. \textbf{A model that cannot tell \textit{"good, not bad"} from \textit{"bad, not good"} still does well on sentiment}, because long reviews carry plenty of unordered evidence.\par}\vskip 4pt
\end{frame}

\framekicker{Section 14.4.2}
\begin{frame}[shrink=25]{Bigrams: putting a little order back by hand}
{\fontsize{9.4}{12.6}\selectfont\color{fg2}A bigram is a pair of adjacent words treated as one token. It restores \textbf{local} order without any change to the model.\par}\vskip 4pt
\icode{python}{listing --- bigram encoding}{listings/ch14-09.py}
\begin{iband}[signal]
One argument, and the model improves: \textit{"not good"} is now a token in its own right. But the approach \hilite{only scales to a local ordering of a few words}, and it is feature engineering done by hand.
\end{iband}
\end{frame}

\framekicker{Section 14.4.1}
\begin{frame}[shrink=25]{Counts, or presence? A third output mode}
{\fontsize{9.4}{12.6}\selectfont\color{fg2}\inlinecode{multi\_hot} records only whether a word appeared. \inlinecode{TextVectorization} offers two alternatives that weight it.\par}\vskip 4pt
\begin{itable}{>{\raggedright\arraybackslash}p{0.2256\textwidth}>{\raggedright\arraybackslash}p{0.7144\textwidth}}
\thead{\inlinecode{output\_mode}} & \thead{What each position holds} \\ \midrule
\inlinecode{"multi\_hot"} & \textbf{1 if the word appears}, 0 otherwise. \\
\inlinecode{"count"} & \textbf{How many times} it appears. \\
\inlinecode{"tf\_idf"} & Its count, \textbf{discounted by how common the word is across the whole corpus} --- so \textit{the} contributes almost nothing and a distinctive word contributes a lot. \\
\end{itable}
\begin{iband}[signal]
TF-IDF is decades older than deep learning and still competitive on short texts. It is worth trying before \hilite{anything with a recurrent layer in it}
\end{iband}
\end{frame}

\coursesection{04}{Sequence models}{Stop engineering order. Show the model the raw sequence.}

\framekicker{Section 14.5}
\begin{frame}[shrink=25]{The argument for moving on}
{\fontsize{9.4}{12.6}\selectfont\color{fg2}The bigram result showed that \textbf{sequence information matters}. But it was obtained by manually engineering features, and that approach has a low ceiling.\par}\vskip 4pt
\begin{iband}[signal]
As so often in deep learning: rather than building the features yourself, \textbf{expose the model to the raw word sequence and let it learn the positional dependencies} --- \hilite{the same argument chapter 1 made about feature engineering in general}
\end{iband}
\end{frame}

\framekicker{Section 14.5}
\begin{frame}[shrink=25]{One wrinkle: batches have to be rectangular}
{\fontsize{9.4}{12.6}\selectfont\color{fg2}Reviews have different lengths, but a batch must be a rectangular tensor so the computation parallelises on a GPU.\par}\vskip 4pt
\icode{python}{padding and truncation}{listings/ch14-10.py}
\begin{iband}[amber]
Padding introduces a new problem the model has to be told about: \textbf{those zeros are not words}. Getting this wrong means \hilite{the model spends capacity learning about padding}
\end{iband}
\end{frame}

\framekicker{Section 14.5.1}
\begin{frame}[shrink=25]{A recurrent model on the raw sequence}
{\fontsize{9.4}{12.6}\selectfont\color{fg2}Chapter 13's machinery applies directly --- a bidirectional LSTM, which here \textbf{does} help, because both directions of a sentence carry meaning.\par}\vskip 4pt
\icode{python}{a sequence model}{listings/ch14-11.py}
\begin{iband}[signal]
\textbf{\inlinecode{mask\_zero=True}} is the answer to the padding problem: it tells every downstream layer to \hilite{skip the padded positions entirely}
\end{iband}
\end{frame}

\framekicker{Section 14.5.1}
\begin{frame}[shrink=25]{What the sequence model buys, and what it costs}
\begin{minipage}[t]{0.494\textwidth}
\begin{icardgood}
\cardh{What it buys}
Word order is \textbf{learned rather than engineered}, so dependencies longer than a bigram become reachable.
\end{icardgood}
\end{minipage}%
\hfill
\begin{minipage}[t]{0.494\textwidth}
\begin{icardwarn}
\cardh{What it costs}
Far slower to train than a bag-of-words model, and it needs \textbf{more data} to justify the extra parameters.
\end{icardwarn}
\end{minipage}%
\begin{iband}[signal]
The book's standing advice applies: \textbf{start with the cheap model}. If bag-of-words is close to the sequence model on your data, \hilite{the sequence model is not earning its cost}
\end{iband}
\end{frame}

\coursesection{05}{Word embeddings}{One-hot encoding makes an assumption that is plainly false.}

\framekicker{Section 14.5.2}
\begin{frame}[shrink=25]{What one-hot encoding quietly assumes}
{\fontsize{9.4}{12.6}\selectfont\color{fg2}Choosing one-hot is a \textbf{feature engineering decision}: it injects an assumption about the structure of the feature space. One-hot vectors are all \textbf{orthogonal to one another} --- the tokens are assumed independent.\par}\vskip 4pt
\begin{iband}[rose]
For words that assumption is \textbf{clearly wrong}. \textit{"Movie"} and \textit{"film"} are interchangeable in most sentences, so their vectors \hilite{should be the same vector, or close to it} --- not at right angles.
\end{iband}
\end{frame}

\framekicker{Section 14.5.2}
\begin{frame}[shrink=25]{What an embedding is instead}
\begin{center}
\includegraphics[width=0.92\linewidth,height=0.60\textheight,keepaspectratio]{figs/ch14-embed.tex.pdf}
\end{center}
\vskip 3pt{\color{ink3}\fontsize{7.4}{9.4}\selectfont Figure 14.3 --- sparse, hardcoded, high-dimensional versus dense, learned, and lower-dimensional.\par}
{\fontsize{9.4}{12.6}\selectfont\color{fg2}The goal: \textbf{the geometric relationship between two word vectors should reflect the semantic relationship between the words}. Related words close, unrelated words far apart.\par}\vskip 4pt
\end{frame}

\framekicker{Section 14.5.2}
\begin{frame}[shrink=25]{Directions in the space mean things}
{\fontsize{9.4}{12.6}\selectfont\color{fg2}Embeddings are not only dense but \textbf{structured}, and the structure is learned. The book's toy example places four words on a plane: \textit{cat}, \textit{dog}, \textit{wolf}, \textit{tiger}.\par}\vskip 4pt
\begin{ibullets}
\item The same vector takes you from \textbf{cat $\rightarrow$ tiger} and from \textbf{dog $\rightarrow$ wolf} --- a \textit{from pet to wild animal} direction.
\item Another takes you from \textbf{dog $\rightarrow$ cat} and \textbf{wolf $\rightarrow$ tiger} --- a \textit{from canine to feline} direction.
\item In real embedding spaces the classic examples are \textbf{gender} and \textbf{plural} vectors: add \textit{female} to \textit{king} and you get \hilite{\textit{queen}}.
\end{ibullets}
\begin{iband}[signal]
Real word-embedding spaces typically feature \textbf{thousands} of such interpretable directions.
\end{iband}
\end{frame}

\framekicker{Section 14.5.3}
\begin{frame}[shrink=25]{Is there one ideal embedding space?}
\begin{iband}[signal]
The book's answer is careful: \textbf{possibly --- but we have yet to compute anything of the sort.} What is available in practice is an embedding learned either \hilite{jointly with your task, or on a large corpus beforehand}
\end{iband}
{\fontsize{9.4}{12.6}\selectfont\color{fg2}An \inlinecode{Embedding} layer trained with the classifier learns a space specialised to \textit{this} task. That is often fine --- and when labelled data is scarce, it is not.\par}\vskip 4pt
\end{frame}

\framekicker{Section 14.5.3}
\begin{frame}[shrink=25]{Choosing the embedding dimension}
{\fontsize{9.4}{12.6}\selectfont\color{fg2}The \inlinecode{output\_dim} of an \inlinecode{Embedding} layer is a capacity decision, and the usual rule from chapter 5 applies in both directions.\par}\vskip 4pt
\begin{itable}{>{\raggedright\arraybackslash}p{0.2444\textwidth}>{\raggedright\arraybackslash}p{0.6956\textwidth}}
\thead{Dimension} & \thead{When it fits} \\ \midrule
\textbf{Too small} & An information bottleneck --- the space cannot separate words that mean different things. Chapter 4's 4-unit layer, again. \\
\textbf{256 -- 1,024} & Typical for large vocabularies, and the range the book works in. \\
\textbf{Too large} & More parameters than the data can constrain; the embedding memorises rather than generalises. \\
\end{itable}
\begin{iband}[signal]
For comparison: one-hot over this vocabulary would be \textbf{20,000 dimensions}. An embedding \hilite{packs more information into far fewer}
\end{iband}
\end{frame}

\coursesection{06}{Pretraining}{Where the unlabelled 25,000 reviews finally earn their place.}

\framekicker{Section 14.5.4}
\begin{frame}[shrink=25]{Why pretraining took over NLP}
{\fontsize{9.4}{12.6}\selectfont\color{fg2}Once you move past small set-based models to sequence models with millions or billions of parameters, text models become \textbf{incredibly data-hungry} --- and you are usually limited by \textbf{how many labelled examples you can find}.\par}\vskip 4pt
\begin{iband}[signal]
The idea: \textbf{devise an unsupervised task that needs no labels}. Pretraining data can be text from a similar domain, or even arbitrary text in the right language. It learns general patterns, \hilite{priming the model before you specialise it}
\end{iband}
\end{frame}

\framekicker{Section 14.5.4}
\begin{frame}[shrink=25]{CBOW: guess the word from its neighbours}
\begin{center}
\includegraphics[width=0.92\linewidth,height=0.60\textheight,keepaspectratio]{figs/ch14-cbow.tex.pdf}
\end{center}
\vskip 3pt{\color{ink3}\fontsize{7.4}{9.4}\selectfont Figure 14.6 --- Continuous Bag of Words, from Mikolov et al. (2013).\par}
{\fontsize{9.4}{12.6}\selectfont\color{fg2}If the surrounding bag contains \textit{sail}, \textit{wave}, and \textit{mast}, you might guess \textit{boat}. \textbf{No labels are needed} --- the supervision comes from the text itself, which is exactly chapter 1's \hilite{self-supervised learning}\par}\vskip 4pt
\end{frame}

\framekicker{Section 14.5.4}
\begin{frame}[shrink=25]{Where the extra data comes from}
{\fontsize{9.4}{12.6}\selectfont\color{fg2}Remember the \inlinecode{unsup/} directory that was skipped during dataset preparation? It holds \textbf{another 25,000 reviews} --- the same size as the training set.\par}\vskip 4pt
\icode{python}{combining labelled and unlabelled text}{listings/ch14-12.py}
\begin{iband}[signal]
\textbf{Doubling the text costs nothing}, because the CBOW task does not care whether a review is positive or negative --- \hilite{it only needs words in context}
\end{iband}
\end{frame}

\framekicker{Section 14.5.5}
\begin{frame}[shrink=25]{Using the pretrained embedding}
{\fontsize{9.4}{12.6}\selectfont\color{fg2}The pretrained weights are loaded into the classifier's \inlinecode{Embedding} layer, which then either stays frozen or is fine-tuned --- the same two options as chapter 8.\par}\vskip 4pt
\icode{python}{loading pretrained weights}{listings/ch14-13.py}
{\fontsize{9.4}{12.6}\selectfont\color{fg2}This is \textbf{transfer learning for text}, structurally identical to feature extraction from a pretrained ConvNet --- \hilite{and it is the pattern the whole of chapter 15 builds on}\par}\vskip 4pt
\end{frame}

\framekicker{Sections 14.4 -- 14.5}
\begin{frame}[shrink=25]{The four models, on one benchmark}
\begin{center}
\includegraphics[width=0.92\linewidth,height=0.60\textheight,keepaspectratio]{figs/ch14-ladder.tex.pdf}
\end{center}
\vskip 3pt{\color{ink3}\fontsize{7.4}{9.4}\selectfont Each step adds either more order information or more outside knowledge.\par}
\begin{iband}[signal]
And the practical caveat: \textbf{the bag-of-words model is close behind, trains in seconds, and is trivial to explain}. On short texts, or with little data, \hilite{it is frequently the right answer}
\end{iband}
\end{frame}

\framekicker{Practice}
\begin{frame}[shrink=25]{Four ways text pipelines go wrong}
\begin{minipage}[t]{0.494\textwidth}
\begin{icardbad}
\cardh{Padding treated as content}
Forgetting \inlinecode{mask\_zero=True}. The model spends capacity modelling zeros, and long reviews are penalised against short ones.
\end{icardbad}
\end{minipage}%
\hfill
\begin{minipage}[t]{0.494\textwidth}
\begin{icardbad}
\cardh{adapt() on the wrong data}
Calling \inlinecode{adapt()} on the full dataset builds the vocabulary from \textbf{test text} --- an information leak in the sense of chapter 5.
\end{icardbad}
\end{minipage}%


\vskip 4pt

\begin{minipage}[t]{0.494\textwidth}
\begin{icardwarn}
\cardh{Truncation that cuts the answer}
\inlinecode{output\_sequence\_length} shorter than the reviews means the ending --- often where sentiment resolves --- is discarded.
\end{icardwarn}
\end{minipage}%
\hfill
\begin{minipage}[t]{0.494\textwidth}
\begin{icardwarn}
\cardh{A tokenizer trained on the wrong language}
A vocabulary built on English text will decompose other languages into near-characters, inflating sequence length enormously.
\end{icardwarn}
\end{minipage}%
\end{frame}
\note{The second one is the subtle one and the most common in submitted work: adapt() is learning from data, so it obeys the same split discipline as training does.}

\framekicker{Summary}
\begin{frame}[shrink=25]{What has to survive this chapter}
\begin{isteps}
\item \textbf{Natural language is messy by construction} --- usage precedes rules, which is why rule-based NLP failed for forty years.
\item \textbf{Tokenization is compression.} Byte-pair encoding gets short sequences and a small vocabulary at once.
\item \textbf{Set or sequence} is the question every NLP architecture answers.
\item \textbf{Bag-of-words is a strong, cheap baseline.} Bigrams add local order for one argument.
\item \textbf{One-hot assumes words are independent}, which is false. Embeddings learn a space where geometry reflects meaning.
\item \textbf{Pretraining on unlabelled text} removes the labelled-data bottleneck.
\end{isteps}
\vskip 4pt
\reslink{NOTEBOOK}{04\_pretraining\_an\_embedding.ipynb}{../../course-slides/notebooks/ch14/04\_pretraining\_an\_embedding.ipynb}
\reslink{NEXT}{Chapter 15 --- Language models and the Transformer}{../ch15/index.html}
\vskip 2pt
\end{frame}

\end{document}
